\cotefile{157-2}
\scanpage{131}
\mspage{76}
% archive sheet number: 130

\diagram{\begin{tikzcd}[ampersand replacement=\&,column sep=0.45em,row sep=0.68em,font=\scriptsize]
\&\&\&\underline{\mathrm{Fibhot}}\,S
 \arrow[dd,hook,"\unsure{\delta|d}"]
 \arrow[ddddlll,dashed,bend right=12,"\Phi_0"']
 \arrow[ddddrrr,dashed,bend left=12,"\Psi_0"]\\
\&\&\underline{\mathrm{Liss}}_0\,S \arrow[ur,hook] \arrow[ddd,hook]\&\&
 \underline{\mathrm{Prop}}_0\,S \arrow[ul,hook] \arrow[ddd,hook]\\
\&\underline{\mathrm{Fib}}_0\,S \arrow[ur,hook] \arrow[ddd,hook]\&
 \&\mathrm{Cat}_{/S} \arrow[d,"C"] \arrow[ddddlll,dashed,bend right=9,"\Phi"'] \arrow[ddddrrr,dashed,bend left=9,"\Psi"]\&\&
 \underline{\mathrm{Cofib}}_0\,S \arrow[ul,hook] \arrow[ddd,hook]\\
\&\&\&\underline{\mathrm{Parf}}\,S \arrow[uul,hook] \arrow[uur,hook]
 \arrow[ull,dotted,leftarrow,"\deleted{\unsure{\triangleleft}}"'{pos=0.94}] \arrow[urr,dotted,leftarrow,"\deleted{\unsure{\triangleright}}"{pos=0.94}]
 \arrow[ddd,dashed,"\mathrm{id}"]\\
\underline{\mathrm{Fibsc}}_0\,S \arrow[uur] \arrow[dd,hook]\&\&
 \underline{\mathrm{Liss}}\,S \arrow[uur,hook]\&\&
 \underline{\mathrm{Prop}}\,S \arrow[uul,hook]\&\&
 \underline{\mathrm{Cofibsc}}_0\,S \arrow[uul] \arrow[dd,hook]\\
\&\underline{\mathrm{Fib}}\,S \arrow[ur,hook]\&\&\&\&
 \underline{\mathrm{Cofib}}\,S \arrow[ul,hook]\\
\underline{\mathrm{Fibsc}}\,S \arrow[ur]\&\&\&
 \underline{\mathrm{Parf}}\,S \arrow[uul,dashed,hook] \arrow[uur,dashed,hook]\&\&\&
 \underline{\mathrm{Cofibsc}}\,S \arrow[ul] \arrow[ul,phantom,"\deleted{\unsure{\downarrow}}"']
\end{tikzcd}}

J'ai \unsure{tracé|dessiné} en \inserted{pointillés} verts les cinq foncteurs $\Phi,\Psi,\Phi_0,\Psi_0,C$ qui \deleted{\unsure{par passage au dérivé}} \inserted{par passage aux localisés} donnent des quasi-inverses pour toutes les équivalences de catégories énoncées dans \textcircled{I} \textcircled{II} du présent \unsure{Scholie}. \deleted{\unsure{De}} Pour préciser, le foncteur

\[
\Phi:\mathrm{Cat}_{/S}\longrightarrow\underline{\mathrm{Fibsc}}\,S
\]

associe à tout $X$ de $\mathrm{Cat}_{/S}$ la catégorie scindée sur $S$

\[
s\longmapsto {}_s\backslash X
\]

sur $S$. Ce foncteur, et ses restrictions à $\underline{\mathrm{Liss}}\,S$, $\underline{\mathrm{Fib}}\,S$ et, plus généralement, à toute sous-catégorie pleine $\unsure{\mathcal{F}|F}$ de $\mathrm{Cat}_{/S}$ contenant $\underline{\mathrm{Fib}}\,S$, \deleted{établit} est \unsure{compatible} avec ces localiseurs $W_S$ d'une part (sur $\mathcal{F}$ et $\underline{\mathrm{Fibsc}}\,S$), $W_S^d$ d'autre part, et passe \inserted{donc} aux localisés correspondants, pour fournir le quasi-inverse \inserted{cherché} de

\[
(\underline{\mathrm{Fibsc}}\,S)\Sigma^{-1}\longrightarrow\mathcal{F}\,\Sigma^{-1}
\qquad(\Sigma=W_S\text{ ou }W_S^d).
\]

On peut aussi, pour tte sous-catégorie pleine $\mathcal{F}'$ intermédiaire

\[
\underline{\mathrm{Fib}}_S\subset\mathcal{F}'\subset\mathcal{F}\subset\mathrm{Cat}_{/S}
\]

regarder $\Phi$ comme un foncteur $\mathcal{F}\longrightarrow\mathcal{F}'$ (plus précisément, prendre le composé $\mathcal{F}\xrightarrow{\Phi|\mathcal{F}}\underline{\mathrm{Fibsc}}\,S\longrightarrow\underline{\mathrm{Fib}}\,S\hookrightarrow\mathcal{F}'$) soit $\Phi_{\mathcal{F}',\mathcal{F}}$, et \deleted{\unsure{de ce qui}} ce qui vient d'être dit résulte que les deux foncteurs

\diagram{\begin{tikzcd}[ampersand replacement=\&,column sep=large]
\mathcal{F}'\arrow[r,hook]\&\mathcal{F}\arrow[l,bend left=16,"\Phi_{\mathcal{F}',\mathcal{F}}"]
\end{tikzcd}}
